\documentclass{beamer}
\usepackage[utf8]{inputenc}
\usepackage[spanish]{babel}
\usepackage{listings}
\usepackage{xcolor}
\usepackage{tikz}
\usepackage{tcolorbox}
\usepackage{fontawesome5}

% Configuración de colores - Gruvbox Light con marrón principal
\definecolor{bg0}{HTML}{fbf1c7}
\definecolor{bg1}{HTML}{ebdbb2}
\definecolor{bg2}{HTML}{d5c4a1}
\definecolor{bg3}{HTML}{bdae93}
\definecolor{fg0}{HTML}{282828}
\definecolor{fg1}{HTML}{3c3836}
\definecolor{fg2}{HTML}{504945}
\definecolor{aqua}{HTML}{689d6a}
\definecolor{aquabright}{HTML}{427b58}
\definecolor{blue}{HTML}{458588}
\definecolor{red}{HTML}{cc241d}
\definecolor{orange}{HTML}{d65d0e}
\definecolor{yellow}{HTML}{d79921}
\definecolor{green}{HTML}{98971a}
\definecolor{purple}{HTML}{b16286}
\definecolor{brown}{HTML}{AB886D}
\definecolor{brownbright}{HTML}{8B6849}

% Colores para código
\definecolor{codegreen}{HTML}{98971a}
\definecolor{codegray}{HTML}{7c6f64}
\definecolor{codepurple}{HTML}{b16286}
\definecolor{backcolour}{HTML}{f9f5d7}

% Configuración de código
\lstdefinestyle{pythonstyle}{
    backgroundcolor=\color{backcolour},   
    commentstyle=\color{codegreen},
    keywordstyle=\color{brownbright}\bfseries,
    numberstyle=\tiny\color{codegray},
    stringstyle=\color{codepurple},
    basicstyle=\ttfamily\scriptsize\color{fg0},
    breakatwhitespace=false,         
    breaklines=true,                 
    captionpos=b,                    
    keepspaces=true,                 
    numbers=left,                    
    numbersep=2pt,                  
    showspaces=false,                
    showstringspaces=false,
    showtabs=false,                  
    tabsize=2,
    language=Python,
    frame=single,
    rulecolor=\color{brown},
    xleftmargin=6pt,
    xrightmargin=6pt,
    inputencoding=utf8,
    extendedchars=true,
    literate=
        {á}{{\'a}}1 {é}{{\'e}}1 {í}{{\'i}}1 {ó}{{\'o}}1 {ú}{{\'u}}1
        {Á}{{\'A}}1 {É}{{\'E}}1 {Í}{{\'I}}1 {Ó}{{\'O}}1 {Ú}{{\'U}}1
        {ñ}{{\~n}}1 {Ñ}{{\~N}}1
        {¿}{{?`}}1 {¡}{{!`}}1
}

\lstset{style=pythonstyle}

% Tema personalizado con marrón principal
\usetheme{Madrid}
\setbeamercolor{structure}{fg=brown}
\setbeamercolor{palette primary}{bg=brown,fg=bg0}
\setbeamercolor{palette secondary}{bg=brownbright,fg=bg0}
\setbeamercolor{palette tertiary}{bg=fg1,fg=bg0}
\setbeamercolor{palette quaternary}{bg=brown,fg=bg0}
\setbeamercolor{block title}{bg=brown,fg=bg0}
\setbeamercolor{block body}{bg=bg1,fg=fg0}
\setbeamercolor{frametitle}{bg=brown,fg=bg0}
\setbeamercolor{title}{fg=fg0}
\setbeamercolor{subtitle}{fg=fg1}
\setbeamercolor{normal text}{fg=fg0,bg=bg0}
\setbeamercolor{item}{fg=brown}
\setbeamercolor{section in toc}{fg=brown}

% Cajas personalizadas con marrón
\newtcolorbox{conceptbox}[1]{
    colback=bg1,
    colframe=brown,
    fonttitle=\bfseries\small,
    title=#1,
    coltitle=fg0,
    left=3pt,
    right=3pt,
    top=3pt,
    bottom=3pt,
    boxsep=2pt
}

\newtcolorbox{notebox}{
    colback=yellow!15,
    colframe=orange,
    leftrule=2mm,
    rightrule=0mm,
    toprule=0mm,
    bottomrule=0mm,
    arc=0mm,
    left=3pt,
    right=3pt,
    top=2pt,
    bottom=2pt
}

\title{\textbf{Introducción a la Programación con Python}}
\subtitle{Semana 9: Conceptos Avanzados}
\author{$\nabla$Nablabs}
\date{}
\begin{document}

\begin{frame}
\titlepage
\end{frame}

\begin{frame}
\frametitle{Contenido}
\tableofcontents
\end{frame}

\section{Conjuntos (sets)}

\begin{frame}{Conjuntos: Sets \faObjectGroup}
\begin{conceptbox}{¿Qué son los sets?}
\small
\begin{itemize}
    \item[\faCircle] Colección \textbf{desordenada} de elementos únicos
    \item[\faCheckDouble] No permite duplicados (los elimina automáticamente)
    \item[\faRandom] Sin índice ni orden garantizado
    \item[\faRocket] Operaciones muy rápidas (membership, unión, intersección)
    \item[\faEdit] Mutable (se puede modificar)
\end{itemize}
\end{conceptbox}

\vspace{0.2cm}
\begin{center}
\large\textcolor{brown}{\faLightbulb\ ¡Perfecto para eliminar duplicados!}
\end{center}
\end{frame}

\begin{frame}[fragile]{Crear sets \faPlus}
\begin{lstlisting}
# Formas de crear sets
numeros = {1, 2, 3, 4, 5}
letras = set(['a', 'b', 'c', 'a'])  # Elimina 'a' duplicada

# Set vacío (¡OJO! {} crea dict, no set)
vacio = set()

# Desde string (elimina duplicados)
caracteres = set("hola")

print(numeros)
print(letras)
print(caracteres)
\end{lstlisting}

\vspace{-0.2cm}
\begin{tcolorbox}[colback=green!10,colframe=green,title={\small\faArrowCircleRight \ Salida},coltitle=fg0,left=3pt,right=3pt,top=2pt,bottom=2pt,boxsep=1pt]
\small\texttt{\{1, 2, 3, 4, 5\}\\
\{'a', 'b', 'c'\}\\
\{'h', 'o', 'l', 'a'\}}
\end{tcolorbox}

\vspace{-0.1cm}
\begin{notebox}
\small\faExclamationTriangle\ El orden de los elementos puede variar
\end{notebox}
\end{frame}

\begin{frame}[fragile]{Operaciones con sets \faTools}
\begin{lstlisting}
frutas = {'manzana', 'pera', 'uva'}

# Agregar
frutas.add('kiwi')
print(f"Agregar: {frutas}")

# Eliminar (error si no existe)
frutas.remove('pera')
print(f"Remove: {frutas}")

# Eliminar (sin error si no existe)
frutas.discard('banana')

# Limpiar
frutas_copia = frutas.copy()
frutas_copia.clear()
print(f"Clear: {frutas_copia}")
\end{lstlisting}

\vspace{-0.2cm}
\begin{tcolorbox}[colback=green!10,colframe=green,title={\small\faArrowCircleRight \ Salida},coltitle=fg0,left=3pt,right=3pt,top=2pt,bottom=2pt,boxsep=1pt]
\small\texttt{Agregar: \{'manzana', 'pera', 'uva', 'kiwi'\}\\
Remove: \{'manzana', 'uva', 'kiwi'\}\\
Clear: set()}
\end{tcolorbox}
\end{frame}

\begin{frame}[fragile]{Operaciones matemáticas de sets \faCalculator}
\begin{lstlisting}
a = {1, 2, 3, 4}
b = {3, 4, 5, 6}

# Unión (todos los elementos)
print(f"Unión: {a | b}")
print(f"Union: {a.union(b)}")

# Intersección (elementos comunes)
print(f"Intersección: {a & b}")
print(f"Intersection: {a.intersection(b)}")

# Diferencia (en a pero no en b)
print(f"Diferencia: {a - b}")

# Diferencia simétrica (no comunes)
print(f"Dif. simétrica: {a ^ b}")
\end{lstlisting}
\end{frame}

\begin{frame}[fragile]{Operaciones matemáticas: Salida \faArrowCircleRight}
\begin{tcolorbox}[colback=green!10,colframe=green,title={\small\faArrowCircleRight \ Salida},coltitle=fg0,left=3pt,right=3pt,top=3pt,bottom=3pt,boxsep=2pt]
\small\texttt{Unión: \{1, 2, 3, 4, 5, 6\}\\
Union: \{1, 2, 3, 4, 5, 6\}\\
Intersección: \{3, 4\}\\
Intersection: \{3, 4\}\\
Diferencia: \{1, 2\}\\
Dif. simétrica: \{1, 2, 5, 6\}}
\end{tcolorbox}

\vspace{0.3cm}
\begin{notebox}
\small\faLightbulb\ Los operadores | \& - \^{} son más legibles que los métodos
\end{notebox}
\end{frame}

\begin{frame}[fragile]{Sets: Caso de uso práctico \faCheckCircle}
\begin{lstlisting}
# Eliminar duplicados de una lista
numeros = [1, 2, 2, 3, 4, 4, 4, 5, 1, 2]
unicos = list(set(numeros))
print(f"Original: {numeros}")
print(f"Sin duplicados: {unicos}")

# Encontrar elementos únicos en dos listas
lista1 = ['python', 'java', 'c++', 'python']
lista2 = ['javascript', 'python', 'ruby']

comunes = set(lista1) & set(lista2)
solo_lista1 = set(lista1) - set(lista2)

print(f"Comunes: {comunes}")
print(f"Solo en lista1: {solo_lista1}")
\end{lstlisting}
\end{frame}

\section{Diccionarios avanzados}

\begin{frame}[fragile]{Diccionarios avanzados \faBook}
\begin{conceptbox}{Métodos útiles}
Más allá de lo básico: métodos avanzados para trabajar con diccionarios
\end{conceptbox}

\begin{lstlisting}
persona = {'nombre': 'Ana', 'edad': 25, 'ciudad': 'Madrid'}

# get() con valor por defecto
print(persona.get('pais', 'No especificado'))

# setdefault() - obtiene o crea con valor default
persona.setdefault('profesion', 'Estudiante')
print(persona)

# update() - actualizar múltiples valores
persona.update({'edad': 26, 'email': 'ana@ejemplo.com'})
print(persona)
\end{lstlisting}
\end{frame}

\begin{frame}[fragile]{Diccionarios: Salida \faArrowCircleRight}
\begin{tcolorbox}[colback=green!10,colframe=green,title={\small\faArrowCircleRight \ Salida},coltitle=fg0,left=3pt,right=3pt,top=2pt,bottom=2pt,boxsep=1pt]
\tiny\texttt{No especificado\\
\{'nombre': 'Ana', 'edad': 25, 'ciudad': 'Madrid', 'profesion': 'Estudiante'\}\\
\{'nombre': 'Ana', 'edad': 26, 'ciudad': 'Madrid', 'profesion': 'Estudiante', 'email': 'ana@ejemplo.com'\}}
\end{tcolorbox}
\end{frame}

\begin{frame}[fragile]{Diccionarios anidados \faLayerGroup}
\begin{lstlisting}
# Diccionarios dentro de diccionarios
empresa = {
    'empleados': {
        'E001': {'nombre': 'Ana', 'salario': 50000},
        'E002': {'nombre': 'Luis', 'salario': 60000}
    },
    'departamentos': {
        'IT': ['E001'],
        'Ventas': ['E002']
    }
}

# Acceder a datos anidados
print(empresa['empleados']['E001']['nombre'])
print(empresa['departamentos']['IT'])

# Iterar sobre diccionario anidado
for id_emp, datos in empresa['empleados'].items():
    print(f"{id_emp}: {datos['nombre']} - ${datos['salario']}")
\end{lstlisting}
\end{frame}

\begin{frame}[fragile]{Diccionarios anidados: Salida \faArrowCircleRight}
\begin{tcolorbox}[colback=green!10,colframe=green,title={\small\faArrowCircleRight \ Salida},coltitle=fg0,left=3pt,right=3pt,top=3pt,bottom=3pt,boxsep=2pt]
\small\texttt{Ana\\
['E001']\\
E001: Ana - \$50000\\
E002: Luis - \$60000}
\end{tcolorbox}
\end{frame}

\section{Funciones lambda}

\begin{frame}[fragile]{Funciones lambda \faLightbulb}
\begin{conceptbox}{¿Qué son?}
Funciones \textbf{anónimas} de una sola expresión. Sintaxis compacta para funciones simples.
\end{conceptbox}

\textbf{Sintaxis:}
\begin{lstlisting}
lambda argumentos: expresion
\end{lstlisting}

\vspace{0.2cm}
\textbf{Comparación:}
\begin{lstlisting}
# Función tradicional
def sumar(a, b):
    return a + b

# Función lambda equivalente
sumar_lambda = lambda a, b: a + b

print(sumar(3, 5))
print(sumar_lambda(3, 5))
\end{lstlisting}

\vspace{-0.2cm}
\begin{tcolorbox}[colback=green!10,colframe=green,title={\small\faArrowCircleRight \ Salida},coltitle=fg0,left=3pt,right=3pt,top=2pt,bottom=2pt,boxsep=1pt]
\small\texttt{8\\
8}
\end{tcolorbox}
\end{frame}

\begin{frame}[fragile]{Lambda con map(), filter(), sorted() \faFilter}
\begin{lstlisting}
numeros = [1, 2, 3, 4, 5]

# map() - aplicar función a cada elemento
cuadrados = list(map(lambda x: x**2, numeros))
print(f"Cuadrados: {cuadrados}")

# filter() - filtrar elementos
pares = list(filter(lambda x: x % 2 == 0, numeros))
print(f"Pares: {pares}")

# sorted() con key personalizada
palabras = ['python', 'es', 'genial', 'y', 'poderoso']
ordenadas = sorted(palabras, key=lambda x: len(x))
print(f"Por longitud: {ordenadas}")

personas = [{'nombre': 'Ana', 'edad': 25}, {'nombre': 'Luis', 'edad': 20}]
por_edad = sorted(personas, key=lambda p: p['edad'])
print(f"Por edad: {por_edad}")
\end{lstlisting}
\end{frame}

\begin{frame}[fragile]{Lambda: Salida \faArrowCircleRight}
\begin{tcolorbox}[colback=green!10,colframe=green,title={\small\faArrowCircleRight \ Salida},coltitle=fg0,left=3pt,right=3pt,top=2pt,bottom=2pt,boxsep=1pt]
\tiny\texttt{Cuadrados: [1, 4, 9, 16, 25]\\
Pares: [2, 4]\\
Por longitud: ['y', 'es', 'genial', 'python', 'poderoso']\\
Por edad: [\{'nombre': 'Luis', 'edad': 20\}, \{'nombre': 'Ana', 'edad': 25\}]}
\end{tcolorbox}

\vspace{0.3cm}
\begin{notebox}
\small\faInfoCircle\ Lambda es útil para funciones simples y temporales
\end{notebox}
\end{frame}

\section{List comprehensions}

\begin{frame}[fragile]{List comprehensions \faListUl}
\begin{conceptbox}{¿Qué son?}
Forma concisa y elegante de crear listas a partir de iterables. Más rápidas y legibles que bucles tradicionales.
\end{conceptbox}

\textbf{Sintaxis:}
\begin{lstlisting}
[expresion for elemento in iterable if condicion]
\end{lstlisting}

\vspace{0.2cm}
\textbf{Comparación:}
\begin{lstlisting}
# Forma tradicional
cuadrados = []
for x in range(5):
    cuadrados.append(x**2)

# List comprehension
cuadrados = [x**2 for x in range(5)]

print(cuadrados)
\end{lstlisting}

\vspace{-0.2cm}
\begin{tcolorbox}[colback=green!10,colframe=green,title={\small\faArrowCircleRight \ Salida},coltitle=fg0,left=3pt,right=3pt,top=2pt,bottom=2pt,boxsep=1pt]
\small\texttt{[0, 1, 4, 9, 16]}
\end{tcolorbox}
\end{frame}

\begin{frame}[fragile]{List comprehensions con condiciones \faFilter}
\begin{lstlisting}
# Números pares
pares = [x for x in range(10) if x % 2 == 0]
print(f"Pares: {pares}")

# Con transformación y condición
cuadrados_impares = [x**2 for x in range(10) if x % 2 != 0]
print(f"Cuadrados impares: {cuadrados_impares}")

# Múltiples condiciones
divisibles = [x for x in range(20) if x % 2 == 0 if x % 3 == 0]
print(f"Divisibles por 2 y 3: {divisibles}")

# If-else (diferente sintaxis)
resultado = ['par' if x % 2 == 0 else 'impar' for x in range(5)]
print(f"Par/Impar: {resultado}")
\end{lstlisting}
\end{frame}

\begin{frame}[fragile]{List comprehensions: Salida \faArrowCircleRight}
\begin{tcolorbox}[colback=green!10,colframe=green,title={\small\faArrowCircleRight \ Salida},coltitle=fg0,left=3pt,right=3pt,top=3pt,bottom=3pt,boxsep=2pt]
\small\texttt{Pares: [0, 2, 4, 6, 8]\\
Cuadrados impares: [1, 9, 25, 49, 81]\\
Divisibles por 2 y 3: [0, 6, 12, 18]\\
Par/Impar: ['par', 'impar', 'par', 'impar', 'par']}
\end{tcolorbox}
\end{frame}

\begin{frame}[fragile]{List comprehensions anidadas \faLayerGroup}
\begin{lstlisting}
# Matriz (lista de listas)
matriz = [[1, 2, 3], [4, 5, 6], [7, 8, 9]]

# Aplanar matriz
plana = [num for fila in matriz for num in fila]
print(f"Aplanada: {plana}")

# Crear matriz con comprehension
matriz_nueva = [[i*j for j in range(3)] for i in range(3)]
print(f"Matriz nueva: {matriz_nueva}")

# Combinar dos listas
colores = ['rojo', 'verde', 'azul']
objetos = ['coche', 'casa']
combinaciones = [f"{color} {objeto}" 
                 for color in colores 
                 for objeto in objetos]
print(f"Combinaciones: {combinaciones}")
\end{lstlisting}
\end{frame}

\begin{frame}[fragile]{List comprehensions anidadas: Salida \faArrowCircleRight}
\begin{tcolorbox}[colback=green!10,colframe=green,title={\small\faArrowCircleRight \ Salida},coltitle=fg0,left=3pt,right=3pt,top=2pt,bottom=2pt,boxsep=1pt]
\tiny\texttt{Aplanada: [1, 2, 3, 4, 5, 6, 7, 8, 9]\\
Matriz nueva: [[0, 0, 0], [0, 1, 2], [0, 2, 4]]\\
Combinaciones: ['rojo coche', 'rojo casa', 'verde coche', 'verde casa', 'azul coche', 'azul casa']}
\end{tcolorbox}
\end{frame}

\section{Dictionary comprehensions}

\begin{frame}[fragile]{Dictionary comprehensions \faBook}
\begin{conceptbox}{¿Qué son?}
Similar a list comprehensions pero para crear diccionarios de forma concisa.
\end{conceptbox}

\textbf{Sintaxis:}
\begin{lstlisting}
{clave: valor for elemento in iterable if condicion}
\end{lstlisting}

\vspace{0.2cm}
\begin{lstlisting}
# Crear diccionario de cuadrados
cuadrados = {x: x**2 for x in range(5)}
print(cuadrados)

# Desde dos listas
claves = ['nombre', 'edad', 'ciudad']
valores = ['Ana', 25, 'Madrid']
persona = {k: v for k, v in zip(claves, valores)}
print(persona)
\end{lstlisting}

\vspace{-0.2cm}
\begin{tcolorbox}[colback=green!10,colframe=green,title={\small\faArrowCircleRight \ Salida},coltitle=fg0,left=3pt,right=3pt,top=2pt,bottom=2pt,boxsep=1pt]
\small\texttt{\{0: 0, 1: 1, 2: 4, 3: 9, 4: 16\}\\
\{'nombre': 'Ana', 'edad': 25, 'ciudad': 'Madrid'\}}
\end{tcolorbox}
\end{frame}

\begin{frame}[fragile]{Dict comprehensions con condiciones \faFilter}
\begin{lstlisting}
# Filtrar diccionario
notas = {'Ana': 85, 'Luis': 92, 'María': 78, 'Pedro': 95}

# Solo aprobados (>= 80)
aprobados = {nombre: nota for nombre, nota in notas.items() 
             if nota >= 80}
print(f"Aprobados: {aprobados}")

# Transformar valores
incrementadas = {nombre: nota + 5 for nombre, nota in notas.items()}
print(f"Incrementadas: {incrementadas}")

# Invertir diccionario
invertido = {v: k for k, v in notas.items()}
print(f"Invertido: {invertido}")
\end{lstlisting}
\end{frame}

\begin{frame}[fragile]{Dict comprehensions: Salida \faArrowCircleRight}
\begin{tcolorbox}[colback=green!10,colframe=green,title={\small\faArrowCircleRight \ Salida},coltitle=fg0,left=3pt,right=3pt,top=2pt,bottom=2pt,boxsep=1pt]
\small\texttt{Aprobados: \{'Ana': 85, 'Luis': 92, 'Pedro': 95\}\\
Incrementadas: \{'Ana': 90, 'Luis': 97, 'María': 83, 'Pedro': 100\}\\
Invertido: \{85: 'Ana', 92: 'Luis', 78: 'María', 95: 'Pedro'\}}
\end{tcolorbox}
\end{frame}

\section{*args y **kwargs}

\begin{frame}[fragile]{Argumentos variables: *args \faEllipsisH}
\begin{conceptbox}{¿Qué es *args?}
Permite pasar un número \textbf{variable} de argumentos posicionales a una función. Se reciben como tupla.
\end{conceptbox}

\begin{lstlisting}
def sumar(*args):
    """Suma cualquier cantidad de números"""
    total = sum(args)
    print(f"Argumentos: {args}")
    return total

print(sumar(1, 2, 3))
print(sumar(10, 20, 30, 40, 50))

def info(nombre, *otros_datos):
    print(f"Nombre: {nombre}")
    print(f"Otros datos: {otros_datos}")

info("Ana", 25, "Madrid", "Ingeniera")
\end{lstlisting}
\end{frame}

\begin{frame}[fragile]{*args: Salida \faArrowCircleRight}
\begin{tcolorbox}[colback=green!10,colframe=green,title={\small\faArrowCircleRight \ Salida},coltitle=fg0,left=3pt,right=3pt,top=3pt,bottom=3pt,boxsep=2pt]
\small\texttt{Argumentos: (1, 2, 3)\\
6\\
Argumentos: (10, 20, 30, 40, 50)\\
150\\
Nombre: Ana\\
Otros datos: (25, 'Madrid', 'Ingeniera')}
\end{tcolorbox}

\vspace{0.3cm}
\begin{notebox}
\small\faInfoCircle\ El nombre 'args' es convención, puedes usar cualquier nombre
\end{notebox}
\end{frame}

\begin{frame}[fragile]{Argumentos variables: **kwargs \faKey}
\begin{conceptbox}{¿Qué es **kwargs?}
Permite pasar un número \textbf{variable} de argumentos con nombre (keyword). Se reciben como diccionario.
\end{conceptbox}

\begin{lstlisting}
def crear_perfil(**kwargs):
    """Crea perfil con datos variables"""
    print("Perfil creado:")
    for clave, valor in kwargs.items():
        print(f"  {clave}: {valor}")
    return kwargs

crear_perfil(nombre="Ana", edad=25, ciudad="Madrid")
print()
crear_perfil(nombre="Luis", profesion="Ingeniero", email="luis@ejemplo.com")
\end{lstlisting}
\end{frame}

\begin{frame}[fragile]{**kwargs: Salida \faArrowCircleRight}
\begin{tcolorbox}[colback=green!10,colframe=green,title={\small\faArrowCircleRight \ Salida},coltitle=fg0,left=3pt,right=3pt,top=3pt,bottom=3pt,boxsep=2pt]
\small\texttt{Perfil creado:\\
\ \ nombre: Ana\\
\ \ edad: 25\\
\ \ ciudad: Madrid\\
\ \\
Perfil creado:\\
\ \ nombre: Luis\\
\ \ profesion: Ingeniero\\
\ \ email: luis@ejemplo.com}
\end{tcolorbox}
\end{frame}

\begin{frame}[fragile]{Combinando *args y **kwargs \faPuzzlePiece}
\begin{lstlisting}
def funcion_completa(arg1, arg2, *args, kwarg1="default", **kwargs):
    print(f"arg1: {arg1}")
    print(f"arg2: {arg2}")
    print(f"args: {args}")
    print(f"kwarg1: {kwarg1}")
    print(f"kwargs: {kwargs}")

funcion_completa(1, 2, 3, 4, kwarg1="personalizado", 
                 extra1="valor1", extra2="valor2")
\end{lstlisting}

\vspace{-0.2cm}
\begin{tcolorbox}[colback=green!10,colframe=green,title={\small\faArrowCircleRight \ Salida},coltitle=fg0,left=3pt,right=3pt,top=2pt,bottom=2pt,boxsep=1pt]
\tiny\texttt{arg1: 1\\
arg2: 2\\
args: (3, 4)\\
kwarg1: personalizado\\
kwargs: \{'extra1': 'valor1', 'extra2': 'valor2'\}}
\end{tcolorbox}
\end{frame}

\section{Generadores e iteradores}

\begin{frame}[fragile]{Generadores \faRecycle}
\begin{conceptbox}{¿Qué son?}
Funciones que generan valores \textbf{uno a la vez} usando \texttt{yield}. Más eficientes en memoria que listas.
\end{conceptbox}

\begin{lstlisting}
# Función generadora
def contador(maximo):
    n = 0
    while n < maximo:
        yield n  # Retorna valor y pausa
        n += 1

# Usar generador
gen = contador(5)
print(type(gen))

for numero in gen:
    print(numero, end=' ')
\end{lstlisting}

\vspace{-0.2cm}
\begin{tcolorbox}[colback=green!10,colframe=green,title={\small\faArrowCircleRight \ Salida},coltitle=fg0,left=3pt,right=3pt,top=2pt,bottom=2pt,boxsep=1pt]
\small\texttt{<class 'generator'>\\
0 1 2 3 4}
\end{tcolorbox}
\end{frame}

\begin{frame}[fragile]{Generadores vs Listas \faBalanceScale}
\begin{lstlisting}
import sys

# Lista (toda en memoria)
lista = [x**2 for x in range(1000)]
print(f"Tamaño lista: {sys.getsizeof(lista)} bytes")

# Generador (valores bajo demanda)
generador = (x**2 for x in range(1000))
print(f"Tamaño generador: {sys.getsizeof(generador)} bytes")

# Usar generador
primeros_5 = [next(generador) for _ in range(5)]
print(f"Primeros 5: {primeros_5}")
\end{lstlisting}

\vspace{-0.2cm}
\begin{tcolorbox}[colback=green!10,colframe=green,title={\small\faArrowCircleRight \ Salida},coltitle=fg0,left=3pt,right=3pt,top=2pt,bottom=2pt,boxsep=1pt]
\small\texttt{Tamaño lista: 8856 bytes\\
Tamaño generador: 104 bytes\\
Primeros 5: [0, 1, 4, 9, 16]}
\end{tcolorbox}

\vspace{-0.1cm}
\begin{notebox}
\small\faLightbulb\ Generadores son ideales para grandes volúmenes de datos
\end{notebox}
\end{frame}

\begin{frame}[fragile]{Generador: Fibonacci \faInfinity}
\begin{lstlisting}
def fibonacci(limite):
    """Genera números de Fibonacci hasta un límite"""
    a, b = 0, 1
    while a < limite:
        yield a
        a, b = b, a + b

# Generar Fibonacci hasta 100
for num in fibonacci(100):
    print(num, end=' ')
print()

# Convertir a lista
fib_lista = list(fibonacci(50))
print(f"\nFibonacci < 50: {fib_lista}")
\end{lstlisting}

\vspace{-0.2cm}
\begin{tcolorbox}[colback=green!10,colframe=green,title={\small\faArrowCircleRight \ Salida},coltitle=fg0,left=3pt,right=3pt,top=2pt,bottom=2pt,boxsep=1pt]
\small\texttt{0 1 1 2 3 5 8 13 21 34 55 89\\
\ \\
Fibonacci < 50: [0, 1, 1, 2, 3, 5, 8, 13, 21, 34]}
\end{tcolorbox}
\end{frame}

\section{Decoradores básicos}

\begin{frame}[fragile]{Decoradores \faMagic}
\begin{conceptbox}{¿Qué son?}
Funciones que modifican o extienden el comportamiento de otras funciones sin cambiar su código.
\end{conceptbox}

\begin{lstlisting}
def mi_decorador(funcion):
    def wrapper():
        print("Antes de la función")
        funcion()
        print("Después de la función")
    return wrapper

@mi_decorador
def saludar():
    print("¡Hola!")

saludar()
\end{lstlisting}

\vspace{-0.2cm}
\begin{tcolorbox}[colback=green!10,colframe=green,title={\small\faArrowCircleRight \ Salida},coltitle=fg0,left=3pt,right=3pt,top=2pt,bottom=2pt,boxsep=1pt]
\small\texttt{Antes de la función\\
¡Hola!\\
Después de la función}
\end{tcolorbox}
\end{frame}

\begin{frame}[fragile]{Decorador con argumentos \faWrench}
\begin{lstlisting}
def decorador_con_args(funcion):
    def wrapper(*args, **kwargs):
        print(f"Llamando a {funcion.__name__}")
        resultado = funcion(*args, **kwargs)
        print(f"Resultado: {resultado}")
        return resultado
    return wrapper

@decorador_con_args
def multiplicar(a, b):
    return a * b

@decorador_con_args
def sumar(a, b, c):
    return a + b + c

multiplicar(3, 4)
print()
sumar(1, 2, 3)
\end{lstlisting}
\end{frame}

\begin{frame}[fragile]{Decorador con argumentos: Salida \faArrowCircleRight}
\begin{tcolorbox}[colback=green!10,colframe=green,title={\small\faArrowCircleRight \ Salida},coltitle=fg0,left=3pt,right=3pt,top=3pt,bottom=3pt,boxsep=2pt]
\small\texttt{Llamando a multiplicar\\
Resultado: 12\\
\ \\
Llamando a sumar\\
Resultado: 6}
\end{tcolorbox}
\end{frame}

\begin{frame}[fragile]{Decorador práctico: Medir tiempo \faStopwatch}
\begin{lstlisting}
import time

def medir_tiempo(funcion):
    def wrapper(*args, **kwargs):
        inicio = time.time()
        resultado = funcion(*args, **kwargs)
        fin = time.time()
        print(f"{funcion.__name__} tardó {fin - inicio:.4f} segundos")
        return resultado
    return wrapper

@medir_tiempo
def calcular_suma(n):
    return sum(range(n))

@medir_tiempo
def calcular_cuadrados(n):
    return [i**2 for i in range(n)]

calcular_suma(1000000)
calcular_cuadrados(10000)
\end{lstlisting}
\end{frame}

\section{Manejo de APIs simples}

\begin{frame}[fragile]{Consumir APIs con requests \faCloud}
\begin{conceptbox}{¿Qué es una API?}
Interfaz que permite comunicación entre aplicaciones. REST API usa HTTP para intercambiar datos (generalmente JSON).
\end{conceptbox}

\textbf{Instalar requests:}
\begin{lstlisting}[language=bash, numbers=none]
pip install requests
\end{lstlisting}

\vspace{0.2cm}
\textbf{Ejemplo básico:}
\begin{lstlisting}
import requests

# GET request a una API pública
url = "https://jsonplaceholder.typicode.com/posts/1"
respuesta = requests.get(url)

print(f"Status: {respuesta.status_code}")
print(f"Tipo: {type(respuesta.json())}")
print(f"Datos: {respuesta.json()}")
\end{lstlisting}
\end{frame}

\begin{frame}[fragile]{API: Salida ejemplo \faArrowCircleRight}
\begin{tcolorbox}[colback=green!10,colframe=green,title={\small\faArrowCircleRight \ Salida},coltitle=fg0,left=3pt,right=3pt,top=2pt,bottom=2pt,boxsep=1pt]
\tiny\texttt{Status: 200\\
Tipo: <class 'dict'>\\
Datos: \{'userId': 1, 'id': 1, 'title': 'sunt aut...', 'body': 'quia et...'\}}
\end{tcolorbox}

\vspace{0.3cm}
\begin{notebox}
\small\faInfoCircle\ Status 200 significa éxito, 404 no encontrado, 500 error del servidor
\end{notebox}
\end{frame}

\begin{frame}[fragile]{API: Obtener múltiples recursos \faList}
\begin{lstlisting}
import requests

# Obtener lista de posts
url = "https://jsonplaceholder.typicode.com/posts"
respuesta = requests.get(url)

if respuesta.status_code == 200:
    posts = respuesta.json()
    print(f"Total de posts: {len(posts)}")
    
    # Mostrar primeros 3
    for post in posts[:3]:
        print(f"\nID: {post['id']}")
        print(f"Título: {post['title']}")
        print(f"Usuario: {post['userId']}")
else:
    print(f"Error: {respuesta.status_code}")
\end{lstlisting}
\end{frame}

\begin{frame}[fragile]{API: POST request \faUpload}
\begin{lstlisting}
import requests

url = "https://jsonplaceholder.typicode.com/posts"

# Datos a enviar
nuevo_post = {
    "title": "Mi nuevo post",
    "body": "Este es el contenido del post",
    "userId": 1
}

# Enviar POST
respuesta = requests.post(url, json=nuevo_post)

print(f"Status: {respuesta.status_code}")
print(f"Respuesta: {respuesta.json()}")
\end{lstlisting}

\vspace{-0.2cm}
\begin{tcolorbox}[colback=green!10,colframe=green,title={\small\faArrowCircleRight \ Salida},coltitle=fg0,left=3pt,right=3pt,top=2pt,bottom=2pt,boxsep=1pt]
\tiny\texttt{Status: 201\\
Respuesta: \{'title': 'Mi nuevo post', 'body': '...', 'userId': 1, 'id': 101\}}
\end{tcolorbox}
\end{frame}

\section{Entornos virtuales}

\begin{frame}[fragile]{Entornos virtuales \faBox}
\begin{conceptbox}{¿Qué son?}
Espacios aislados para cada proyecto con sus propias dependencias y versiones de paquetes.
\end{conceptbox}

\textbf{¿Por qué usarlos?}
\small
\begin{itemize}
    \item[\faCheckCircle] Evitar conflictos entre proyectos
    \item[\faCheckCircle] Mantener dependencias específicas por proyecto
    \item[\faCheckCircle] Facilitar colaboración y despliegue
    \item[\faCheckCircle] No contaminar instalación global de Python
\end{itemize}
\end{frame}

\begin{frame}[fragile]{Crear y usar entornos virtuales \faTerminal}
\textbf{Crear entorno virtual:}
\begin{lstlisting}[language=bash, numbers=none]
# Python 3
python -m venv mi_entorno

# O con virtualenv
virtualenv mi_entorno
\end{lstlisting}

\vspace{0.2cm}
\textbf{Activar entorno:}
\begin{lstlisting}[language=bash, numbers=none]
# Linux/Mac
source mi_entorno/bin/activate

# Windows
mi_entorno\Scripts\activate
\end{lstlisting}

\vspace{0.2cm}
\textbf{Desactivar:}
\begin{lstlisting}[language=bash, numbers=none]
deactivate
\end{lstlisting}
\end{frame}

\begin{frame}[fragile]{Gestión de dependencias \faListAlt}
\textbf{Instalar paquetes:}
\begin{lstlisting}[language=bash, numbers=none]
# Con entorno activado
pip install requests pandas numpy
\end{lstlisting}

\vspace{0.2cm}
\textbf{Guardar dependencias:}
\begin{lstlisting}[language=bash, numbers=none]
pip freeze > requirements.txt
\end{lstlisting}

\vspace{0.2cm}
\textbf{Instalar desde requirements.txt:}
\begin{lstlisting}[language=bash, numbers=none]
pip install -r requirements.txt
\end{lstlisting}

\vspace{0.2cm}
\begin{notebox}
\small\faLightbulb\ requirements.txt facilita compartir y replicar entornos
\end{notebox}
\end{frame}

\begin{frame}[fragile]{Ejemplo de requirements.txt \faFileAlt}
\begin{tcolorbox}[colback=bg1,colframe=brown,left=3pt,right=3pt,top=2pt,bottom=2pt]
\small\texttt{requests==2.31.0\\
pandas==2.0.3\\
numpy==1.24.3\\
matplotlib==3.7.2\\
pytest==7.4.0}
\end{tcolorbox}

\vspace{0.3cm}
\textbf{Buenas prácticas:}
\small
\begin{itemize}
    \item[\faCheck] Un entorno virtual por proyecto
    \item[\faCheck] Incluir requirements.txt en control de versiones
    \item[\faCheck] Documentar versión de Python requerida
    \item[\faCheck] Actualizar requirements.txt regularmente
    \item[\faCheck] No incluir carpeta del entorno en git (usar .gitignore)
\end{itemize}
\end{frame}

\section{Proyecto integrador}

\begin{frame}{Proyecto: Analizador de datos desde API \faChartLine}
\begin{conceptbox}{Objetivo}
Crear una aplicación que consuma una API, procese datos y genere estadísticas usando conceptos avanzados.
\end{conceptbox}

\textbf{Funcionalidades:}
\small
\begin{itemize}
    \item[\faCloud] Consumir API REST pública
    \item[\faFilter] Filtrar y transformar datos con comprehensions
    \item[\faChartBar] Calcular estadísticas usando lambda y funciones
    \item[\faSave] Guardar resultados en archivo
    \item[\faStopwatch] Medir tiempos de ejecución con decoradores
    \item[\faRecycle] Usar generadores para eficiencia
\end{itemize}

\vspace{0.2cm}
\begin{center}
\textcolor{brown}{\faLightbulb\ Integra todos los conceptos de la semana}
\end{center}
\end{frame}

\begin{frame}[fragile]{Proyecto: Decoradores y utilidades \faTools}
\begin{lstlisting}
import time
import json

def medir_tiempo(func):
    """Decorador para medir tiempo de ejecución"""
    def wrapper(*args, **kwargs):
        inicio = time.time()
        resultado = func(*args, **kwargs)
        fin = time.time()
        print(f"⏱ {func.__name__} tardó {fin - inicio:.2f}s")
        return resultado
    return wrapper

def guardar_json(func):
    """Decorador para guardar resultados en JSON"""
    def wrapper(*args, **kwargs):
        resultado = func(*args, **kwargs)
        nombre_archivo = f"{func.__name__}_resultado.json"
        with open(nombre_archivo, 'w') as f:
            json.dump(resultado, f, indent=2)
        print(f"💾 Guardado en {nombre_archivo}")
        return resultado
    return wrapper
\end{lstlisting}
\end{frame}

\begin{frame}[fragile]{Proyecto: Clase principal \faCode}
\begin{lstlisting}
import requests

class AnalizadorAPI:
    def __init__(self, base_url):
        self.base_url = base_url
        self.datos = []
    
    @medir_tiempo
    def obtener_datos(self, endpoint):
        """Obtiene datos de la API"""
        url = f"{self.base_url}/{endpoint}"
        respuesta = requests.get(url)
        
        if respuesta.status_code == 200:
            self.datos = respuesta.json()
            print(f"✓ Obtenidos {len(self.datos)} registros")
            return self.datos
        else:
            print(f"✗ Error: {respuesta.status_code}")
            return []
\end{lstlisting}

\small(Continuación en siguiente slide...)
\end{frame}

\begin{frame}[fragile]{Proyecto: Filtros y análisis \faFilter}
\begin{lstlisting}
    def filtrar_por_usuario(self, user_id):
        """Filtra datos por ID de usuario"""
        return [item for item in self.datos 
                if item.get('userId') == user_id]
    
    def generar_estadisticas(self):
        """Genera estadísticas de los datos"""
        if not self.datos:
            return {}
        
        # Usar set para usuarios únicos
        usuarios = {item['userId'] for item in self.datos}
        
        # Dict comprehension para contar por usuario
        posts_por_usuario = {
            user_id: len([p for p in self.datos if p['userId'] == user_id])
            for user_id in usuarios
        }
        
        return {
            'total_registros': len(self.datos),
            'usuarios_unicos': len(usuarios),
            'posts_por_usuario': posts_por_usuario
        }
\end{lstlisting}
\end{frame}

\begin{frame}[fragile]{Proyecto: Generador de títulos \faRecycle}
\begin{lstlisting}
    def generar_titulos(self, limite=None):
        """Generador que produce títulos uno a uno"""
        contador = 0
        for item in self.datos:
            if limite and contador >= limite:
                break
            yield item.get('title', 'Sin título')
            contador += 1
    
    def buscar_palabra_en_titulos(self, palabra):
        """Busca palabra en títulos usando lambda"""
        # Filter con lambda
        coincidencias = list(filter(
            lambda item: palabra.lower() in item.get('title', '').lower(),
            self.datos
        ))
        return coincidencias
\end{lstlisting}
\end{frame}

\begin{frame}[fragile]{Proyecto: Métodos de ordenamiento \faSort}
\begin{lstlisting}
    def top_usuarios(self, n=5):
        """Top N usuarios con más posts"""
        stats = self.generar_estadisticas()
        posts = stats['posts_por_usuario']
        
        # Sorted con lambda para ordenar por valor
        top = sorted(posts.items(), 
                    key=lambda x: x[1], 
                    reverse=True)[:n]
        
        return dict(top)
    
    @guardar_json
    def generar_reporte(self):
        """Genera reporte completo"""
        return {
            'estadisticas': self.generar_estadisticas(),
            'top_usuarios': self.top_usuarios(3),
            'total_posts': len(self.datos)
        }
\end{lstlisting}
\end{frame}

\begin{frame}[fragile]{Proyecto: Uso del analizador \faTerminal}
\begin{lstlisting}
# Crear analizador
analizador = AnalizadorAPI("https://jsonplaceholder.typicode.com")

# Obtener datos
analizador.obtener_datos("posts")

# Estadísticas generales
stats = analizador.generar_estadisticas()
print(f"\n📊 Estadísticas:")
print(f"  Total registros: {stats['total_registros']}")
print(f"  Usuarios únicos: {stats['usuarios_unicos']}")

# Top usuarios
print(f"\n🏆 Top 3 usuarios:")
top = analizador.top_usuarios(3)
for user_id, count in top.items():
    print(f"  Usuario {user_id}: {count} posts")
\end{lstlisting}

\small(Continuación en siguiente slide...)
\end{frame}

\begin{frame}[fragile]{Proyecto: Uso - Generadores y filtros \faPlay}
\begin{lstlisting}
# Usar generador de títulos
print(f"\n📝 Primeros 5 títulos:")
for i, titulo in enumerate(analizador.generar_titulos(5), 1):
    print(f"  {i}. {titulo[:50]}...")

# Buscar palabra
print(f"\n🔍 Posts sobre 'qui':")
resultados = analizador.buscar_palabra_en_titulos("qui")
print(f"  Encontrados: {len(resultados)}")
for post in resultados[:3]:
    print(f"  - {post['title'][:40]}...")

# Generar reporte (usa decorador @guardar_json)
print(f"\n📄 Generando reporte...")
analizador.generar_reporte()
\end{lstlisting}
\end{frame}

\begin{frame}[fragile]{Proyecto: Salida esperada \faArrowCircleRight}
\begin{tcolorbox}[colback=green!10,colframe=green,title={\small\faArrowCircleRight \ Salida},coltitle=fg0,left=3pt,right=3pt,top=2pt,bottom=2pt,boxsep=1pt]
\tiny\texttt{⏱ obtener\_datos tardó 0.45s\\
✓ Obtenidos 100 registros\\
\ \\
📊 Estadísticas:\\
\ \ Total registros: 100\\
\ \ Usuarios únicos: 10\\
\ \\
🏆 Top 3 usuarios:\\
\ \ Usuario 1: 10 posts\\
\ \ Usuario 2: 10 posts\\
\ \ Usuario 3: 10 posts\\
\ \\
📝 Primeros 5 títulos:\\
\ \ 1. sunt aut facere repellat provident occaecati...\\
\ \ 2. qui est esse...\\
\ \\
🔍 Posts sobre 'qui':\\
\ \ Encontrados: 15\\
\ \\
📄 Generando reporte...\\
💾 Guardado en generar\_reporte\_resultado.json}
\end{tcolorbox}
\end{frame}

\begin{frame}{Proyecto alternativo: Gestor de tareas avanzado \faClipboardList}
\begin{conceptbox}{Alternativa}
Sistema de gestión de tareas con características avanzadas
\end{conceptbox}

\textbf{Características:}
\small
\begin{itemize}
    \item[\faListUl] CRUD de tareas con comprehensions
    \item[\faFilter] Filtros múltiples con lambda
    \item[\faTag] Sistema de etiquetas con sets
    \item[\faChartBar] Estadísticas y reportes
    \item[\faSave] Persistencia en JSON
    \item[\faStopwatch] Tracking de tiempo con decoradores
    \item[\faRecycle] Generador de tareas pendientes
    \item[\faBell] Recordatorios y notificaciones
\end{itemize}
\end{frame}

\begin{frame}{Proyecto alternativo: Web scraper \faSpider}
\begin{conceptbox}{Otra alternativa}
Extractor de datos de sitios web con análisis
\end{conceptbox}

\textbf{Características:}
\small
\begin{itemize}
    \item[\faGlobe] Scraping con requests y BeautifulSoup
    \item[\faFilter] Limpieza de datos con comprehensions
    \item[\faDatabase] Almacenamiento en estructuras eficientes
    \item[\faChartLine] Análisis estadístico
    \item[\faFileExport] Exportación a múltiples formatos
    \item[\faStopwatch] Monitoreo de rendimiento
    \item[\faRecycle] Procesamiento lazy con generadores
\end{itemize}
\end{frame}

\begin{frame}{Buenas prácticas avanzadas \faCheckDouble}
\begin{conceptbox}{Recomendaciones}
\small
\begin{itemize}
    \item[\faCheck] Usa comprehensions para código más pythonic
    \item[\faCheck] Prefiere generadores para grandes datasets
    \item[\faCheck] Lambda solo para funciones muy simples
    \item[\faCheck] Decoradores para separar concerns (logging, timing, etc.)
    \item[\faCheck] Sets para operaciones de conjunto y eliminar duplicados
    \item[\faCheck] Entornos virtuales siempre, uno por proyecto
    \item[\faCheck] Maneja errores en requests (try-except, status codes)
    \item[\faCheck] Documenta código complejo
    \item[\faCheck] Type hints para código más claro
    \item[\faCheck] No abuses de las características avanzadas
\end{itemize}
\end{conceptbox}
\end{frame}

\begin{frame}[fragile]{Type hints (bonus) \faTag}
\begin{lstlisting}
from typing import List, Dict, Optional, Set

def procesar_numeros(
    numeros: List[int], 
    operacion: str = "suma"
) -> int:
    """Procesa lista de números con operación especificada"""
    if operacion == "suma":
        return sum(numeros)
    elif operacion == "max":
        return max(numeros)
    return 0

def crear_perfil(
    nombre: str, 
    edad: int, 
    hobbies: Optional[Set[str]] = None
) -> Dict[str, any]:
    """Crea diccionario de perfil"""
    return {
        "nombre": nombre,
        "edad": edad,
        "hobbies": hobbies or set()
    }
\end{lstlisting}
\end{frame}

\begin{frame}{Recursos adicionales \faBook}
\begin{conceptbox}{Para profundizar}
\small
\begin{itemize}
    \item[\faGlobe] \textbf{Python Docs}: docs.python.org
    \item[\faGlobe] \textbf{Real Python}: realpython.com
    \item[\faGlobe] \textbf{PEP 8}: Guía de estilo oficial
    \item[\faBook] \textbf{Fluent Python}: Libro avanzado
    \item[\faGithub] \textbf{Awesome Python}: Lista de recursos
    \item[\faPlayCircle] \textbf{Corey Schafer}: YouTube tutorials
    \item[\faCode] \textbf{LeetCode/HackerRank}: Práctica
\end{itemize}
\end{conceptbox}

\vspace{0.2cm}
\begin{center}
\textcolor{brown}{\faRocket\ ¡La práctica hace al maestro!}
\end{center}
\end{frame}

\begin{frame}{Resumen \faListUl}
\begin{conceptbox}{¿Qué aprendimos hoy?}
\small
\begin{itemize}
    \item[\faCheck] Sets: colecciones únicas y operaciones matemáticas
    \item[\faCheck] Diccionarios avanzados y anidados
    \item[\faCheck] Funciones lambda para código conciso
    \item[\faCheck] List y dict comprehensions
    \item[\faCheck] *args y **kwargs para argumentos variables
    \item[\faCheck] Generadores para eficiencia de memoria
    \item[\faCheck] Decoradores para extender funcionalidad
    \item[\faCheck] Consumo de APIs con requests
    \item[\faCheck] Entornos virtuales para aislar proyectos
    \item[\faCheck] Proyecto: Analizador de datos desde API
\end{itemize}
\end{conceptbox}

\vspace{0.3cm}
\begin{center}
\large\textcolor{brown}{\faRocket\ ¡Python avanzado desbloqueado!}
\end{center}
\end{frame}

\begin{frame}
\begin{center}
\Huge\textcolor{brown}{\faQuestionCircle}\\
\vspace{0.5cm}
\Huge ¡Preguntas!\\
\vspace{1cm}
\Large\textcolor{brownbright}{Gracias por su atención}
\end{center}
\end{frame}

\end{document}
